\chapter{Integrating published and publishable studies}\label{ch:article-integration}
\section{State the common question}
An article-based thesis needs an argument that connects the studies. Explain
what each study contributes, how they address a common question and why their
combination produces a contribution greater than a sequence of summaries.
Retain a substantive introduction and conclusion.

\section{Record publication and contribution status}
For each incorporated study, record its title, authors, publication status,
version and permission position. Explain your contribution to collaborative
work precisely. Confirm how earlier publications count towards the doctorate
under the regulations applicable to your candidature.

\begin{table}[htbp]
\centering\small
\caption{Editable article-integration record}\label{tab:articles}
\begin{tabularx}{\textwidth}{@{}p{0.23\textwidth}X@{}}
\toprule
Field & Information to supply \\
\midrule
Study & Title, research question and thesis chapter \\
Status & Submitted, accepted or published; exact version and date \\
Contribution & Research, analysis and writing attributable to each author \\
Rights & Applicable licence, permission or retained right \\
Connection & Relationship to the common question and other chapters \\
\bottomrule
\end{tabularx}
\end{table}

\section{Resolve differences between the studies}
Explain changes in terminology, methods, legal cut-off dates and conclusions.
Identify overlaps and cross-references. A published article's editorial format
does not itself determine the thesis format. Confirm which changes and versions
your university and the relevant rights holder permit.

\LawPrompt{Insert your studies in separate chapter files and replace this
demonstration with a substantive account of their relationship.}
